\documentclass[12pt,a4paper]{article}
\usepackage[utf8]{inputenc}
\usepackage[T1]{fontenc}
\usepackage[brazilian]{babel}
\usepackage{geometry}
\geometry{top=3cm,bottom=2cm,left=3cm,right=2cm}
\setlength{\headheight}{13.6pt}
\addtolength{\topmargin}{-1.6pt}
\usepackage{setspace}
\onehalfspacing
\usepackage{indentfirst}
\usepackage{hyperref}
\hypersetup{colorlinks=true,linkcolor=blue,citecolor=blue,urlcolor=blue}
\usepackage{bookmark}
\usepackage{fancyhdr}
\pagestyle{fancy}
\fancyhf{}
\fancyhead[L]{\small Planejamento de Pesquisa}
\fancyhead[R]{\small \thepage}
\renewcommand{\headrulewidth}{0.4pt}
\usepackage{times}
\usepackage{enumitem}
\usepackage{longtable}
\usepackage{colortbl}
\usepackage{xcolor}
\usepackage{graphicx}
\usepackage{titlesec}

\titleformat{\section}{\large\bfseries}{\thesection}{1em}{}
\titleformat{\subsection}{\normalsize\bfseries}{\thesubsection}{1em}{}
\titleformat{\subsubsection}{\normalsize\bfseries\itshape}{\thesubsubsection}{1em}{}

\begin{document}

\begin{titlepage}
\begin{center}
\vspace*{3cm}
{\Large\textbf{PLANEJAMENTO DE PESQUISA}}\\[1.5cm]
{\LARGE\textbf{A Proteção da Personalidade Humana na Era da Inteligência Artificial: Contribuições da Encíclica Magnifica Humanitas para o Direito Contemporâneo}}\\[2cm]
{\large Documento de Planejamento --- Estrutura, Metodologia e Protocolo de Auditoria}\\[1cm]
\vfill
{\large \today}
\end{center}
\end{titlepage}

\newpage
\tableofcontents
\newpage

% ============================================================
\section{Apresentação da Pesquisa}
% ============================================================

\subsection{Tema e Delimitação}

O presente projeto de pesquisa insere-se na intersecção entre o Direito Civil (direitos da personalidade), o Direito Digital e a Filosofia do Direito, tendo como objeto central a proteção da personalidade humana diante dos desafios suscitados pela inteligência artificial, à luz dos fundamentos antropológicos e éticos propostos pela Encíclica \textit{Magnifica Humanitas}.

\textbf{Delimitação}: A investigação circunscreve-se ao ordenamento jurídico brasileiro contemporâneo, com incursões no direito comparado e na tradição do pensamento social cristão, especificamente no magistério pontifício da Encíclica \textit{Magnifica Humanitas}.

\subsection{Problema de Pesquisa}

Como os fundamentos antropológicos e éticos da Encíclica \textit{Magnifica Humanitas} podem contribuir para a construção de um marco jurídico de proteção da personalidade humana diante das novas tecnologias de inteligência artificial no direito contemporâneo?

\subsection{Hipóteses}

\begin{enumerate}[label=H\arabic*.]
\item A \textit{Magnifica Humanitas} oferece um conceito de dignidade personalista que supera a visão meramente patrimonialista dos direitos da personalidade.
\item Os princípios éticos da Encíclica (transparência, responsabilidade, subsidiariedade) podem ser traduzidos em critérios normativos para regulação da IA.
\item O direito contemporâneo, especialmente o brasileiro, carece de uma fundamentação antropológica robusta para enfrentar os desafios da IA, lacuna que a Encíclica pode preencher.
\item A proteção da personalidade na era digital exige uma reinterpretação dos institutos clássicos do direito civil à luz de uma antropologia integral.
\end{enumerate}

\subsection{Objetivos}

\textbf{Objetivo Geral:}

Analisar as contribuições da Encíclica \textit{Magnifica Humanitas} para a formulação de um regime jurídico de proteção da personalidade humana na era da inteligência artificial.

\textbf{Objetivos Específicos:}

\begin{enumerate}[label=OE\arabic*.]
\item Sistematizar os fundamentos antropológicos e éticos da \textit{Magnifica Humanitas} relevantes para o direito.
\item Identificar os principais desafios jurídicos à proteção da personalidade humana suscitados pela IA (vieses algorítmicos, manipulação de dados, tomada de decisão autônoma, responsabilidade civil).
\item Examinar a adequação dos institutos clássicos dos direitos da personalidade (nome, imagem, honra, privacidade) aos novos contextos tecnológicos.
\item Propor critérios normativos inspirados na \textit{Magnifica Humanitas} para a regulação da IA no ordenamento brasileiro.
\item Estabelecer um protocolo rigoroso de fichamento e auditoria de citações que garanta rastreabilidade e transparência científica.
\end{enumerate}

\subsection{Justificativa}

A relevância científica reside na lacuna doutrinária identificada: há abundante produção sobre direito digital e proteção de dados, mas escassa investigação que integre o magistério pontifício recente ao debate jurídico sobre IA. A relevância social decorre da urgência de fundamentos antropologicamente consistentes para a regulação de tecnologias que afetam a própria compreensão do humano. A relevância acadêmica consiste na proposição de um protocolo metodológico auditável de fichamento de citações como padrão de transparência científica.

% ============================================================
\section{Estrutura Metodológica}
% ============================================================

\subsection{Tipo de Pesquisa}

Pesquisa jurídico-teórica, de natureza qualitativa, com abordagem hermenêutico-crítica.

\subsection{Procedimentos Técnicos}

\begin{itemize}
\item \textbf{Pesquisa bibliográfica}: levantamento sistematizado em bases nacionais e internacionais.
\item \textbf{Pesquisa documental}: análise da Encíclica \textit{Magnifica Humanitas} e de legislação correlata (LGPD, Marco Civil da Internet, PL 4/2025, PL 2.338/2023, AI Act europeu).
\item \textbf{Análise hermenêutica}: interpretação dos textos à luz do referencial teórico adotado.
\end{itemize}

\subsection{Técnicas de Coleta e Análise}

\begin{itemize}
\item Busca estruturada por descritores em bases indexadas.
\item Leitura sistemática com fichamento padronizado (protocolo de auditoria).
\item Análise de conteúdo categorial-temática.
\item Cruzamento de fontes (direito, teologia moral, filosofia, ciência da computação).
\end{itemize}

% ============================================================
\section{Marco Teórico Preliminar}
% ============================================================

\subsection{Eixos Temáticos}

\begin{enumerate}
\item \textbf{Direitos da Personalidade e Dignidade Humana}
    \begin{itemize}
    \item Teoria tridimensional dos direitos da personalidade (Bittar, 2023).
    \item Dignidade como fundamento ontológico (Sarlet, 2024).
    \item Tutela póstuma da personalidade digital (Doneda, 2023; Schreiber, 2024).
    \end{itemize}

\item \textbf{Inteligência Artificial e Desafios Jurídicos}
    \begin{itemize}
    \item Algoritmos, vieses e discriminação (Zuboff, 2021; Pasquale, 2022).
    \item Tomada de decisão autônoma e responsabilidade civil (Falcão, 2024).
    \item Regulação da IA no Brasil e na União Europeia (AI Act, 2024).
    \item Personalidade digital, herança digital e perfis póstumos (Reis, 2024).
    \end{itemize}

\item \textbf{Encíclica Magnifica Humanitas e Pensamento Social Cristão}
    \begin{itemize}
    \item Antropologia personalista e dignidade ontológica.
    \item Princípios: dignidade, subsidiariedade, solidariedade, destino universal dos bens.
    \item Tecnologia a serviço da pessoa humana — crítica ao tecnocracia.
    \item Relação com a tradição: \textit{Rerum Novarum}, \textit{Populorum Progressio}, \textit{Laudato Si'}.
    \end{itemize}
\end{enumerate}

% ============================================================
\section{Sumário Preliminar do Manuscrito}
% ============================================================

\begin{enumerate}[label=CAPÍTULO \arabic*:]
\item \textbf{Introdução}
    \begin{itemize}
    \item Contextualização temática
    \item Problema, hipóteses e objetivos
    \item Justificativa e metodologia
    \item Estrutura do trabalho
    \end{itemize}

\item \textbf{Fundamentos Antropológicos e Éticos da Magnifica Humanitas}
    \begin{itemize}
    \item 2.1. Contexto histórico-doutrinário da Encíclica
    \item 2.2. A pessoa humana como centro: dignidade ontológica e relacional
    \item 2.3. Princípios ético-sociais: subsidiariedade, solidariedade e destino universal
    \item 2.4. Crítica à tecnocracia e a técnica a serviço da pessoa
    \item 2.5. Diálogo com a tradição do magistério social
    \end{itemize}

\item \textbf{Inteligência Artificial e os Desafios à Personalidade Humana}
    \begin{itemize}
    \item 3.1. Ontologia e classificação dos sistemas de IA
    \item 3.2. Vieses algorítmicos, discriminação e dignidade
    \item 3.3. Autonomia algorítmica, opacidade e responsabilização
    \item 3.4. Dados pessoais, perfilização e identidade digital
    \item 3.5. Herança digital e a personalidade póstuma na rede
    \end{itemize}

\item \textbf{A Proteção Jurídica da Personalidade no Direito Brasileiro}
    \begin{itemize}
    \item 4.1. Direitos da personalidade: evolução dogmática e fundamentos
    \item 4.2. Tutela do nome, imagem, honra e privacidade
    \item 4.3. Aplicação aos ambientes digitais: LGPD, Marco Civil, STJ
    \item 4.4. Insuficiências e lacunas do modelo vigente
    \end{itemize}

\item \textbf{Contribuições da Magnifica Humanitas para a Regulação da IA}
    \begin{itemize}
    \item 5.1. Dignidade personalista como critério normativo
    \item 5.2. Subsidiariedade e governança participativa da IA
    \item 5.3. Transparência, explicabilidade e prestação de contas
    \item 5.4. Solidariedade e inclusão digital
    \item 5.5. Proposta de critérios para avaliação ético-jurídica de sistemas de IA
    \end{itemize}

\item \textbf{Conclusão}
    \begin{itemize}
    \item Síntese dos argumentos
    \item Confirmação ou refutação das hipóteses
    \item Contribuições originais e limites da pesquisa
    \item Agenda de investigações futuras
    \end{itemize}

\item \textbf{Referências}
\item \textbf{Anexos}
    \begin{itemize}
    \item Fichamentos completos (auditáveis)
    \item Índice de citações com rastreio
    \end{itemize}
\end{enumerate}

% ============================================================
\section{Protocolo de Fichamentos e Auditoria de Citações}
% ============================================================

\subsection{Justificativa do Protocolo}

Todo trabalho acadêmico rigoroso exige rastreabilidade completa de suas fontes. O presente protocolo estabelece que \textbf{cada citação} do manuscrito deve ser documentada em um fichamento individual contendo:

\begin{itemize}
\item Referência completa (ABNT).
\item Trecho original (idioma de publicação).
\item Tradução livre para o português.
\item Interpretação e justificativa da relevância para o parágrafo.
\item Resenha crítica para auditoria com verificação de:
    \begin{itemize}
    \item Fidelidade ao original.
    \item Respeito ao contexto original.
    \item Ausência de anacronismo ou descontextualização.
    \item Qualidade da tradução.
    \item Nota de confiança (0--10).
    \end{itemize}
\end{itemize}

\subsection{Estrutura do Fichamento Padrão}

\medskip
\noindent\fbox{%
\begin{minipage}{\dimexpr\textwidth-2\fboxsep-2\fboxrule-12pt\relax}
\small
\textbf{Fichamento Nº XXX} --- \textbf{Autor, Ano}\\

\textbf{Citação no trabalho}: [trecho exato]\\

\textbf{Referência completa}: AUTOR. \textit{Título}. Edição. Local: Editora, Ano. p. XX--YY.\\

\textbf{Rodapé}: [Autor, ano, p. XX--YY]\\

\textbf{Trecho original}: "[texto original]"\\

\textbf{Tradução}: "[tradução]"\\

\textbf{Interpretação e justificativa}: [...]\\

\textbf{Resenha crítica}:
\begin{itemize}
\item Fidelidade ao original? [Sim/Não]
\item Contexto respeitado? [Sim/Não]
\item Anacronismo? [Sim/Não]
\item Tradução adequada? [Sim/Não]
\item Nota de confiança: [0--10]
\end{itemize}
\end{minipage}
}

\subsection{Organização dos Arquivos}

Os fichamentos serão armazenados na pasta \texttt{fichamentos/} com a seguinte convenção:

\begin{itemize}
\item \texttt{N001\_Autor\_PalavraChave.md}
\item \texttt{N002\_Autor\_PalavraChave.md}
\item \texttt{N000\_INDICE.md} (registro mestre)
\end{itemize}

\subsection{Critérios de Auditoria}

\begin{enumerate}[label=A\arabic*.]
\item \textbf{Fidelidade}: a citação reproduz exatamente o sentido do original, sem distorção ou recorte tendencioso.
\item \textbf{Contexto}: o trecho citado mantém o significado que possuía na obra original, considerando o argumento do autor.
\item \textbf{Anacronismo}: a citação não é aplicada a contexto histórico, tecnológico ou jurídico incompatível com sua origem.
\item \textbf{Tradução}: a versão em português preserva o conteúdo semântico e as nuances do original.
\item \textbf{Justificativa}: a relevância da citação para o argumento no parágrafo é explicitamente demonstrada.
\end{enumerate}

% ============================================================
\section{Plano de Busca Bibliográfica}
% ============================================================

\subsection{Bases de Dados}

\medskip
{\small
\begin{longtable}{|p{4.5cm}|p{3.5cm}|p{5.5cm}|}
\hline
\rowcolor[gray]{0.9}
\textbf{Base} & \textbf{Tipo} & \textbf{Cobertura} \\
\hline
\endfirsthead
\hline
\rowcolor[gray]{0.9}
\textbf{Base} & \textbf{Tipo} & \textbf{Cobertura} \\
\hline
\endhead
SciELO Brasil & Periódicos & Direito, Ciências Sociais \\
\hline
Portal CAPES & Periódicos & Multidisciplinar \\
\hline
Google Scholar & Geral & Multidisciplinar \\
\hline
SSRN & Preprints & Direito \\
\hline
Scopus & Indexada & Direito, Filosofia \\
\hline
Web of Science & Indexada & Direito, Filosofia \\
\hline
Vatican.va & Documental & Magistério Pontifício \\
\hline
Zenith (OpenCode) & Jurisprudência & STJ, STF \\
\hline
\end{longtable}
}

\subsection{Descritores de Busca}

\medskip
{\small
\begin{longtable}{|p{5.5cm}|p{8cm}|}
\hline
\rowcolor[gray]{0.9}
\textbf{Português} & \textbf{Inglês/Italiano} \\
\hline
\endfirsthead
\hline
\rowcolor[gray]{0.9}
\textbf{Português} & \textbf{Inglês/Italiano} \\
\hline
\endhead
Proteção da personalidade & Personality protection \\
\hline
Direitos da personalidade & Personality rights / Diritti della personalità \\
\hline
Inteligência artificial e direito & Artificial intelligence and law \\
\hline
Magnifica Humanitas & Magnifica Humanitas \\
\hline
Dignidade humana e tecnologia & Human dignity and technology \\
\hline
Herança digital & Digital inheritance \\
\hline
Personalidade póstuma & Post-mortem personality \\
\hline
Responsabilidade civil IA & AI tort liability \\
\hline
Vieses algorítmicos & Algorithmic bias \\
\hline
Encíclica e direito & Encyclical and law \\
\hline
\end{longtable}
}

\subsection{Critérios de Inclusão e Exclusão}

\textbf{Critérios de inclusão}:
\begin{itemize}
\item Publicações em português, inglês, italiano e espanhol.
\item Período: 2000--2026 (com prioridade para 2020--2026 em IA).
\item Relevância temática direta para um ou mais eixos da pesquisa.
\item Qualis A1--B2 (periódicos) ou reconhecimento acadêmico (livros).
\end{itemize}

\textbf{Critérios de exclusão}:
\begin{itemize}
\item Publicações sem revisão por pares (blogs, newsletters não acadêmicas).
\item Obras cujo acesso integral não seja verificável.
\item Duplicatas entre bases.
\end{itemize}

% ============================================================
\section{Cronograma Preliminar}
% ============================================================

\medskip
\begin{longtable}{|p{3cm}|p{12cm}|}
\hline
\rowcolor[gray]{0.9}
\textbf{Fase} & \textbf{Atividades} \\
\hline
\endfirsthead
\hline
\rowcolor[gray]{0.9}
\textbf{Fase} & \textbf{Atividades} \\
\hline
\endhead
Fase 1 & Planejamento --- elaboração do projeto, definição do protocolo \\
\hline
Fase 2 & Levantamento bibliográfico --- buscas sistematizadas e fichamento \\
\hline
Fase 3 & Leitura e análise da Encíclica Magnifica Humanitas \\
\hline
Fase 4 & Redação dos capítulos 1 e 2 \\
\hline
Fase 5 & Redação dos capítulos 3 e 4 \\
\hline
Fase 6 & Redação do capítulo 5 e conclusão \\
\hline
Fase 7 & Revisão, auditoria de citações e formatação \\
\hline
Fase 8 & Depósito e defesa \\
\hline
\end{longtable}

% ============================================================
\section{Referências Iniciais}
% ============================================================

\textbf{Direitos da Personalidade e Dignidade:}

BITTAR, Carlos Alberto. \textit{Os direitos da personalidade}. 9. ed. São Paulo: Saraiva, 2023.

SARLET, Ingo Wolfgang. \textit{Dignidade da pessoa humana e direitos fundamentais}. 12. ed. Porto Alegre: Livraria do Advogado, 2024.

SCHREIBER, Anderson. \textit{Direitos da personalidade}. 5. ed. Rio de Janeiro: Forense, 2024.

DONEDA, Danilo. \textit{Da privacidade à proteção de dados pessoais}. 4. ed. São Paulo: RT, 2023.

\textbf{Inteligência Artificial e Direito:}

ZUBOFF, Shoshana. \textit{A era do capitalismo de vigilância}. Rio de Janeiro: Intrínseca, 2021.

PASQUALE, Frank. \textit{Novas leis da robótica}. São Paulo: Fósforo, 2022.

FALCÃO, Eduardo. \textit{Responsabilidade civil e inteligência artificial}. Rio de Janeiro: Forense, 2024.

REIS, Fernando. \textit{Herança digital: proteção de dados e personalidade póstuma}. São Paulo: Almedina, 2024.

\textbf{Magistério Pontifício:}

FRANCISCO, Papa. \textit{Carta Encíclica Magnifica Humanitas}. Vaticano: Libreria Editrice Vaticana, 2025.

FRANCISCO, Papa. \textit{Laudato Si'}. Vaticano: LEV, 2015.

PAULO VI, Papa. \textit{Populorum Progressio}. Vaticano: LEV, 1967.

LEÃO XIII, Papa. \textit{Rerum Novarum}. Vaticano: LEV, 1891.

\textbf{Legislação e Regulação:}

BRASIL. \textit{Lei nº 13.709, de 14 de agosto de 2018} (LGPD). Brasília, 2018.

BRASIL. \textit{Lei nº 12.965, de 23 de abril de 2014} (Marco Civil da Internet). Brasília, 2014.

UNIÃO EUROPEIA. \textit{Regulamento (UE) 2024/1689} (AI Act). Bruxelas, 2024.

BRASIL. \textit{Projeto de Lei nº 2.338, de 2023} (Marco Legal da IA). Brasília, 2023.

BRASIL. \textit{Projeto de Lei nº 4, de 2025} (Reforma do Código Civil). Brasília, 2025.

\end{document}
